\documentclass{resume} % Use the custom resume.cls style

\usepackage[left=0.4 in,top=0.4in,right=0.4 in,bottom=0.4in]{geometry}

\newcommand{\tab}[1]{\hspace{.2667\textwidth}\rlap{#1}}
\newcommand{\itab}[1]{\hspace{0em}\rlap{#1}}
\name{Чигиринцев Игорь}

\address{+7(708) 772-18-04 \\ Казахстан, Алматы}
\address{Email: \href{mailto:chigirintsevigor@gmail.com}{chigirintsevigor@gmail.com} \\ Telegram: \href{https://t.me/cigorn}{@cigorn} }

\begin{document}

\begin{rSection}{Образование}

{\bf Научная математика, бакалавриат}, КазНУ им Аль-Фараби, 4-й курс. \hfill {2021 - 2025}
\end{rSection}

\begin{rSection}{Опыт работы}

\textbf{Ведущий разработчик} \hfill {Август 2024 — июнь 2026}\\
ЭЙЧАР МЕД
\begin{itemize}
    \item Разработка и поддержка клиентской части веб-приложения.
    \item Тестирование баз данных и серверной части веб-приложения.
    \item Написание и тестирование программного кода, исправление выявленных багов.
    \item Регулярный мониторинг работоспособности и производительности программного продукта.
    \item Анализ причин багов с целью их предотвращения; формализация и алгоритмизация поставленных задач.
\end{itemize}

\textbf{Программист-разработчик} \hfill {Июнь 2021 — май 2026}\\
CIN.dev
\begin{itemize}
    \item Fullstack-разработка собственных продуктов компании: платформа объявлений \href{https://tugo.kz}{tugo.kz} и мобильная игра «Sudoku».
    \item Написание и тестирование программного кода, исправление выявленных багов.
\end{itemize}

\end{rSection}

\begin{rSection}{Навыки}
\begin{itemize}
    \item \textit{Языки программирования:} C++, Python3, SQL, Golang, JavaScript, HTML, CSS, Dart.
    \item \textit{Фреймворки:} React, Flutter, Astro, Tailwind.
    \item \textit{Алгоритмы:} классические алгоритмы и структуры данных, цифровая обработка сигналов (FFT, спектральный анализ).
    \item \textit{Базы данных:} PostgreSQL, MS Access, MySQL
    \item \textit{Серверные т.:} Docker-compose, Nginx, Caddy.
    \item \textit{Сетевые т.:} gRPC, API, HTTP, HTTPS.
    \item \textit{Другое:} Unit testing, *nix, Git, Redis, Kafka, интеграция LLM (Gemini API).
\end{itemize}
\end{rSection}

\begin{rSection}{Проекты}
\vspace{-1.25em}
\item \textbf{«Стук» — диагностика неисправностей автомобиля по звуку (\href{https://pro-stuk.com}{pro-stuk.com}).} {Самостоятельная разработка продукта целиком в едином монорепозитории: Android-приложение, бэкенд и сайт. Бэкенд на Go (chi, slog, REST API): собственный DSP-пакет без внешних зависимостей — radix-2 FFT, спектральный центроид, поиск пиков, огибающая и автокорреляция для оценки частоты ударов и оборотов двигателя, покрыт юнит-тестами на синтетических сигналах; мультимодальная интеграция с Gemini API (structured output по JSON-схеме, аудио inline base64) с классификацией пригодности записи, дневными лимитами на устройство, rate-limiting по IP и таймаутами. Мобильное приложение на Flutter (Android): запись WAV PCM16 16 кГц, дерево решений из 44 вопросов и 82 исходов как общий источник данных для приложения и сайта, отчёт с вероятностями причин и светофором срочности, шеринг результата через RepaintBoundary, подписанная release-сборка (Gradle, flutter\_launcher\_icons). Сайт на Astro + Tailwind + React-остров: 170+ статических страниц, pSEO-контур из 150 статей с автоматической перелинковкой, schema.org (Article, FAQPage, BreadcrumbList), sitemap с lastmod; собственная обезличенная аналитика без cookies и сторонних трекеров (просмотры, скачивания приложения, конверсия, фильтрация ботов). Развёрнуто в Docker на VPS за Caddy (авто-TLS) и Cloudflare параллельно с другим проектом на том же сервере.}
\item \textbf{Платформа объявлений «Tugo» (\href{https://tugo.kz}{tugo.kz}).} {Самостоятельная fullstack-разработка и вывод в продакшен платформы объявлений Казахстана (авто, недвижимость, работа, услуги). Бэкенд на Go (chi, pgx/pgxpool, PostgreSQL + PostGIS для геопоиска и карты), JWT-аутентификация, REST API. Фронтенд на React + TypeScript + Vite + Tailwind с мультиязычностью (ru/kz/en) и серверным рендером страниц для поисковых ботов. Реализованы поиск с фильтрами по вертикалям, интерактивная карта на PostGIS, персональные рекомендации, чат, модерация, резюме и вакансии. Полная SEO-оптимизация (локализованные URL, hreflang, sitemap, schema.org) — PageSpeed SEO 100/100. Развёрнуто в Docker (Caddy, авто-TLS) на VPS в Казахстане за Cloudflare (DDoS/WAF); настроены rate-limiting, защита от типовых уязвимостей и автоматические бэкапы.}
\item \textbf{Автоматизированная агропромышленная система.} {Постановка задач команде, разработка общей логики проекта, анализ и оптимизация производительности.}
\item \textbf{Бот реферал.} {Написание алгоритма сообщений бота, построение общей логики, оптимизация производительности. Бот разработан с использованием Python3 и Telegram API (aiogram3) \href{https://t.me/Cin_Snip_Bot}{@Cin\_Snip\_Bot}.}
\item \textbf{Лэндинг учебного центра.} {Разработка структуры сайта, адаптивный дизайн, настройка взаимодействия через HTML, CSS и JavaScript.}
\item \textbf{Мобильное приложение "Watch\&Earn".} {Разработка мобильного приложения под Android на Flutter с использованием архитектурного паттерна Provider для управления состоянием. Реализована интеграция Google AdMob SDK для монетизации приложения через показ рекламы. Организована работа с локальной базой данных с использованием SQLite (sqflite) и безопасное хранение токенов через Flutter Secure Storage. Настроена аутентификация и взаимодействие с SQL базами, которые хостятся на Supabase, с использованием JWT для защиты пользовательских данных. Выполнена адаптивная верстка с использованием Material Design, добавлена оптимизация интерфейса для различных устройств с помощью flutter displaymode. Проведена оптимизация сборки через Gradle, уменьшение размера APK с помощью ProGuard и R8, релизная версия подписана для публикации.}
\item \textbf{Бот парсер.} {Разработка Python-скрипта для Telegram-бота, предназначенного для парсинга цен с веб-сайтов. Используется библиотека requests для отправки HTTP-запросов и BeautifulSoup для парсинга HTML-страниц. Бот обрабатывает как статические, так и динамические сайты, вычисляя среднюю цену по URL и Xpath из предоставленного Excel-файла. Результаты отправляются пользователю в виде отчёта.}
\item \textbf{Мобильная игра "Sudoku".} {Разработка кроссплатформенного мобильного приложения с судоку на Flutter. Реализованы четыре режима сетки (4×4, 6×6, 9×9, 16×16), онлайн-мультиплеер в реальном времени с системой матчмейкинга на основе Supabase Realtime, таблица лидеров, система достижений и друзей. Аутентификация выполнена через Google Sign-In и Supabase Auth с JWT. Монетизация реализована через Google AdMob и внутриигровые покупки (in-app purchase). Добавлена поддержка офлайн-режима с синхронизацией данных, система тем оформления и многоязычный интерфейс. Релизная сборка оптимизирована с помощью Gradle, ProGuard и R8.}
\end{rSection}

\begin{rSection}{Языки}
\begin{itemize}
    \item Русский --- Родной
    \item Английский --- Upper-Intermediate
    \item Казахский --- Разговорный
\end{itemize}
\end{rSection}

\end{document}
